% --------------------------------------------------
% 4. EMPIRICAL STRATEGY
% --------------------------------------------------
\section{Empirical Strategy}
\label{sec:empirical_strategy}

\subsection{Cox Proportional Hazard Model}

The central methodological choice is the Cox proportional hazard (PH) model \parencite{cox1972regression, cleves2010survival}, implemented using the \texttt{lifelines} library for Python \parencite{davidson2023lifelines}. The model is well suited to the data for three reasons. First, it is semi-parametric: it does not require a parametric assumption about the shape of the baseline hazard, which matters because there is no strong prior about whether abandonment risk follows a Weibull or exponential distribution over stages. Second, the Cox model naturally handles right-censoring: riders who complete all 21 stages are correctly accounted for as observations where the event did not occur. Third, the model captures duration dependence through the shared baseline hazard, allowing abandonment risk to evolve across stages without imposing a functional form.

The hazard function for rider $i$ at stage $t$ is:
\begin{equation}
    h(t \mid \mathbf{X}_i) = h_0(t) \exp\!\left(\mathbf{X}_i \boldsymbol{\beta}\right),
\end{equation}
where $h_0(t)$ is the unspecified baseline hazard, $\mathbf{X}_i$ is a vector of covariates, and $\boldsymbol{\beta}$ is the vector of log-hazard coefficients. The hazard ratio $\exp(\beta_k)$ is the key quantity: values above 1 indicate higher abandonment risk, values below 1 indicate a protective effect.

The basic model includes: rider age and previous Grand Tour starts (human capital proxies), historical DNF rate (continuation cost proxy), a GC contender indicator (tournament incentive proxy), team size and team average age (team-resource controls), race dummies for the Tour de France and Giro d'Italia with the Vuelta as baseline (event-value proxies), and a centred year trend.

\subsection{Time-Varying Extension}

A limitation of the basic model is that it assigns fixed covariate values for the entire race, but stage difficulty changes day by day. To address this, a time-varying Cox model adds four stage-type dummies (mountain, hilly, individual time trial, team time trial; flat as baseline) and standardised stage distance. These covariates are entered in the counting-process (start-stop) format, where each row represents one stage interval for one rider \parencite{cleves2010survival}. This permits a direct test of whether stage-level physical costs affect abandonment independently of rider-level characteristics.

\subsection{Proportional Hazards Assumption}

Both models rest on the proportional hazards assumption: the hazard ratio between any two riders is constant over time. This is tested using Schoenfeld residuals \parencite{schoenfeld1982residuals}. Table~\ref{tab:schoenfeld} in the Appendix reports the test statistics: no covariate is significant at the 5\% level, providing no evidence against the assumption. The Schoenfeld residual plots (Figure~\ref{fig:schoenfeld} in the Appendix) confirm visually that there is no systematic trend over time.
